\section{Introduction}\label{Section:Introduction}

Solving multivariate polynomial systems over finite fields is a basic task in computational algebra. Three major viewpoints dominate the area: Gr\"obner basis methods~\cite{sturmfels2002solving,Faugere2002,faugere1999new}, resultant-based elimination~\cite{cox2005using,macaulay1902some,DBLP:conf/issac/KapurSY94}, and characteristic-set methods in the spirit of Wu's method~\cite{Wang1996characteristic}. In contemporary algebraic cryptanalysis, most attention is paid to Gr\"obner basis techniques~\cite{DBLP:journals/iacr/SauerS21,DBLP:journals/tosc/KoschatkoLR24}, while resultant methods---especially the Dixon resultant as a compact multivariate elimination tool---remain much less explored in the algebraic cryptanalysis literature.

These problems are especially relevant for \emph{arithmetization-oriented}~(AO) primitives, whose round functions are deliberately designed to admit compact arithmetic representations in zero-knowledge (ZK) proof systems, MPC, and homomorphic settings. AO primitives differ substantially from classical binary-oriented designs. Instead of working over small binary fields, they are typically defined over large prime fields or extension fields, and their nonlinear layers are optimized for arithmetic circuits. Following this design philosophy, the ZK-friendly AO primitives are commonly organized into three families (Types~I--III)\footnote{See the STAP Zoo taxonomy of ZK-friendly designs: \url{https://stap-zoo.com/zk/}.}:
\begin{itemize}[nosep]
	\item \emph{Type~I (low degree):} MiMC~\cite{DBLP:conf/asiacrypt/Albrecht0R0T16}, GMiMC~\cite{DBLP:conf/esorics/Albrecht0PRRR0S19}, \poseidon~\cite{DBLP:conf/uss/0001KR0S21}, \poseidonnew~\cite{DBLP:conf/africacrypt/GrassiKS23}, \texttt{Neptune}~\cite{DBLP:journals/tosc/GrassiOPS22};
	\item \emph{Type~II (equivalence relations):} \vision/\texttt{Rescue}~\cite{DBLP:journals/tosc/AlyABDS20}, \texttt{XHash12}~\cite{DBLP:journals/iacr/AshurKM23}, \texttt{Anemoi}~\cite{DBLP:conf/crypto/BouvierBCPSVW23}, \texttt{Arion}~\cite{DBLP:journals/corr/abs-2303-04639}, \texttt{Griffin}~\cite{Griffin2023};
	\item \emph{Type~III (lookup tables):} \texttt{Reinforced Concrete}~\cite{ReinforcedConcrete2022}, \texttt{Monolith}~\cite{Monolith2024}, \texttt{Tip5}~\cite{Tip5_2023}.
\end{itemize}
In all three settings, algebraic preimage and collision attacks against these primitives reduce to \emph{constrained-input constrained-output}~\cite{bertoni2007sponge} (CICO) problems, which in turn amount to solving induced polynomial systems whose complexity dominates the security analysis.

\subsubsection{Related works.}
Solving-based algebraic attacks have seen continuous refinement over the last several years. A major line of work~\cite{DBLP:journals/iacr/SauerS21,DBLP:conf/crypto/BariantBLAOPR24,DBLP:journals/tosc/Steiner24,DBLP:journals/tosc/KoschatkoLR24,DBLP:conf/asiacrypt/LiuMM24,DBLP:journals/tosc/GrassiKR25,DBLP:journals/iacr/BakBBBOP25} studies Gr\"obner basis attacks and, in particular, how to obtain better algebraic models or tighter upper bounds for their complexity on structured AO systems. In~\cite{DBLP:journals/tosc/KoschatkoLR24}, the authors refined the complexity analysis by providing more precise estimates for the individual stages of the attack. However, in some cases the upper bounds obtained from such generic Gr\"obner basis analyses remain inaccurate. For example, the \emph{Freelunch} approach~\cite{DBLP:conf/crypto/BariantBLAOPR24} involves leveraging dedicated (weighted) monomial orderings for which a given polynomial system is already a Gr\"obner basis, allowing for better complexity bounds.

Another line of work brings classical elimination theory into cryptanalysis. Yang et al.~\cite{DBLP:conf/asiacrypt/YangZYLT24} introduced iterative pairwise resultant attacks based on the Sylvester resultant, eliminating one variable at a time from two equations. Bariant et al.~\cite{DBLP:conf/crypto/BariantBBHOR25} further improved this paradigm by exploiting structure inside the polynomial system and performing degree reductions during elimination. These are important developments, but they remain rooted in iterated two-polynomial resultants rather than \emph{multipolynomial} resultants, so the elimination order is comparatively rigid.

\paragraph{Dixon resultant.}
The Dixon resultant, introduced by A.C. Dixon in 1909~\cite{Dixon1909}, offers a different determinant-based elimination mechanism: it can eliminate several variables simultaneously from a polynomial system and thus serves as a compact multipolynomial resultant in our setting. Compared with the Macaulay resultant~\cite{macaulay1902some}, the Dixon construction often leads to substantially smaller matrices while preserving the essential elimination information. Kapur et al.~\cite{DBLP:conf/issac/KapurSY94} also explained how to handle the non-square situation via maximal-rank submatrices, which makes the method viable beyond the classical idealized setting. 
Tang and Feng~\cite{cryptoeprint:2005/312} turned this into a solver by keeping one variable as a parameter, and Albrecht~\cite{Albrecht2008} implemented the same workflow against a block-cipher system; our direct method follows it, the additions being degree-aware submatrix selection and a fast polynomial-matrix determinant stage.
Subsequently, numerous works~\cite{lewis2015dixon,Zhao2005extfast,DengQiTangDuan2022,DBLP:journals/jsc/QinZYC22,DBLP:journals/ijcm/QinWTJ17,DBLP:conf/casc/JinaduM22,JinaduMonagan2025} studied both the algebraic and algorithmic aspects of Dixon matrices.

For complexity analysis, the size of the Dixon matrix is the first key parameter. Kapur and Saxena~\cite{DBLP:conf/stoc/KapurS96} gave a Newton-polytope characterization for unmixed systems together with volume-based bounds for multi-homogeneous systems, and Tang and Feng~\cite{cryptoeprint:2005/312} treated the degree-$2$ case. Mixed supports were studied by Chtcherba and Kapur~\cite{ChtcherbaKapur2004} in the bivariate Dixon formulation; we generalize this to an arbitrary number of variables and obtain an explicit ordered mixed-total-degree formula. A second key parameter is the cost of polynomial-matrix determinants themselves, as expansion by minors~\cite{GentlemanJohnson1976}, Bareiss elimination~\cite{HorowitzSahni1975}, Kronecker substitution~\cite{Li2007symbolic}, Hermite normal form~\cite{DBLP:journals/jc/LabahnNZ17}, evaluation-interpolation~\cite{HorowitzSahni1975}, and sparse interpolation~\cite{DBLP:journals/jsc/Huang23} can lead to substantially different complexity estimates, and we accordingly develop a stage-wise complexity model that compares all of them.

\subsubsection*{Our contributions}
This paper studies Dixon resultants as a method for \emph{solving polynomial systems}, with AO primitives as a representative application domain. Our main contributions are as follows.

\begin{enumerate}
	\item We derive refined bounds on Dixon matrix size for mixed systems via lattice-path counting. For uniform-degree systems, the bound specializes to the Fuss--Catalan numbers~\cite{hilton1991catalan,DBLP:journals/dm/Aval08}; for general mixed systems, we obtain an explicit (Hessenberg) determinant formula, which is attained by the monomial support of generic instances.
	
	\item We develop a stage-wise complexity model for Dixon elimination by comparing multiple determinant computation strategies, yielding a more faithful description of polynomial-matrix behavior.
	
	\item We implement an efficient C library \drsolve\footnote{\url{https://github.com/drsolve/drsolve}} supporting multiple determinant routines and a degree-aware submatrix heuristic. Experiments validate the proposed Dixon-matrix bounds and show competitiveness with Magma and msolve on well-determined random systems, with advantages in the low-variable/high-degree/large-field regime, while Magma and msolve remain superior in the high-variable/low-degree/small-field regime.
	
	\item We organize Dixon-based elimination for AO-induced systems into three strategies: direct, iterative, and hybrid reduction. We illustrate them on \poseidon, \vision, and \texttt{XHash12}, respectively, providing concrete algebraic complexity estimates.
\end{enumerate}

\subsubsection{Outline.} 
Section~\ref{Section:Preliminaries} reviews the Dixon construction and the determinant computation models used in the paper. Section~\ref{Section: Complexity} develops the refined Dixon-matrix bounds and the stepwise complexity analysis. Section~\ref{Section: Improvements} discusses implementation-oriented improvements. Section~\ref{Section:Application} presents the three application patterns on \poseidon, \vision, and \xhash.

